% Figure 1: the initial harness H0 and its nine editable surfaces.
% Body only; the `harness' listing style, its colours, and \hzero live in the
% preamble of proposal.tex, matching how Figure_2.tex relies on the tikz styling
% defined there.
\begin{figure}[!tbp]
\centering
\begin{lstlisting}[style=harness]
# shrlm/rlm_harness.py: the nine surfaces Self-Harness may edit (S1 to S9).

MINIMAL_REPL_CONTRACT = """
You are a Recursive Language Model (RLM): a language model whose prompt-related
information lives in a Python REPL that you drive turn by turn.
Available in the REPL:
- `context`: the information related to the prompt (`str` or `list[str]`).
- `llm_query(prompt, model=None) -> str`; `llm_query_batched(prompts, ...)`.
- `rlm_query(...)` / `rlm_query_batched(...)`: recursive RLM sub-calls.
- `SHOW_VARS() -> str`: list every variable currently in the REPL.
- `answer`: dict {"content": "", "ready": False}; set both to submit.
{custom_tools_section}
Write code in ```repl``` blocks, one per turn; the REPL persists. Only
`print(...)` output is shown back. Outputs over ~20K characters are truncated.
"""

def build_repl_contract() -> str:                                    # S1
    return MINIMAL_REPL_CONTRACT

def build_decomposition_instruction() -> str:                        # S2
    return ("Start by probing `context` to understand what you have, "
            "then decide how the task breaks into steps.")

def build_execution_instruction() -> str:                            # S3
    return ("On each turn, run one block of code, print what you need "
            "to see, and build the result up in REPL variables.")

def build_verification_instruction() -> str:                         # S4
    return 'Check your answer before setting `answer["ready"] = True`.'

def build_recovery_instruction() -> str:                             # S5
    return ("If a sub-call errors or returns something you cannot use, "
            "adjust the approach and try again.")

def build_runtime_policy() -> dict[str, Any]:                        # S6
    return {
        "enabled": False,
        "max_prompt_chars": None,
        "max_batch_width": None,
        "max_depth": None,
        "retry_on_syntax_error": None,
        "max_retries": None,
        "validate_sub_output": None,
    }

def build_metadata(                                                  # S7
    stdout: str,
    repl_inventory: dict[str, tuple[str, int]],
    max_character_length: int = DEFAULT_MAX_CHARACTER_LENGTH,
) -> str:
    return default_metadata_builder(stdout, repl_inventory,
                                    max_character_length)

def build_repl_helpers() -> dict[str, Any]:                          # S8
    return {}

def build_sub_repl_helpers() -> dict[str, Any]:                      # S8
    return {}

def build_answer_middleware() -> AnswerMiddlewareFn:                 # S9
    return accept_answer
\end{lstlisting}
\caption{Initial harness \hzero{} and the editable interface used as the starting
point for \sh{}: nine declared surfaces, one builder function each, keyed to phases
of the \rlm{} turn loop. The harness is intentionally kept sparse, with seven of the
nine returning nothing, a disabled policy, or a single generic line, so that the
orchestration strategy a tuned harness carries is something the optimization loop
must recover from its own execution traces. \sid{S6}, \sid{S7} and \sid{S9} required
a runtime injection point in the reference implementation of \citet{zhang2025rlm};
the rest are reached through its existing prompt and tool configuration. A proposal
modifies exactly one surface.}
\label{fig:edit-surfaces}
\end{figure}
